\section{Boundedness of the Individual Position Errors}
\label{sec:bounded}

Section~\ref{sec:avg} showed that only the \emph{average} position error $\bx$ is forced to decay:
it obeys the exponentially stable oscillator \eqref{eq:osc}, so $\bx(t)\to0$. That result proves the
target is fenced (P1), but it says nothing about the \emph{individual} errors
$\txi_i=x_i-x_0$. In principle the average could vanish while each vehicle is far from the target,
arranged so that their centroid cancels exactly. The purpose of this section is to rule that out for
the first controller \eqref{eq:ctrl7}: every individual error $\txi_i(t)$ is bounded. This is a
genuinely different statement from the trichotomy of Conjecture~\ref{conj:dist}, which specifies the
\emph{limit set} of the shape trajectory; here we only establish boundedness of the shape, which is
the weaker fact on which that conjecture implicitly rests.

Throughout we work in the target frame, where the closed loop \eqref{eq:veh}--\eqref{eq:tgt} with
\eqref{eq:ctrl7} reads (cf.\ \eqref{eq:dirn_dyn})
\begin{equation}
  \dot{\txi}_i=\phi_i-k_1\txi_i+\tvell_i,\qquad
  \dot{\tvell}_i=-k_2\txi_i,
  \label{eq:bounded_dyn}
\end{equation}
with $\phi_i=\sum_{j\in\calN_i}\alpha(\norm{x_{ij}})\,x_{ij}/\norm{x_{ij}}$, $x_{ij}=\txi_i-\txi_j$,
and $\alpha$ as in \eqref{eq:alpha} (we extend $\alpha$ by $\alpha(s)=0$ for $s>\mu$, so that the sum
$\sum_{i\ne j}$ below can be taken over all ordered pairs).

\subsection{A pairwise identity for the repulsion}
\label{sec:pairid}

The first structural fact is that the total power of the repulsion, weighted by the individual errors,
collapses to a sum over unordered pairs.

\begin{lemma}[Pairwise repulsion identity]\label{lem:pairid}
Along \eqref{eq:bounded_dyn},
\begin{equation}
  \sum_{i\in\NN}\txi_i^{\mathsf T}\phi_i
  =\frac12\sum_{i\ne j}\alpha(\norm{x_{ij}})\,\norm{x_{ij}} .
  \label{eq:pairid}
\end{equation}
\end{lemma}

\begin{proof}
By definition of $\phi_i$,
\begin{equation}
  \sum_{i\in\NN}\txi_i^{\mathsf T}\phi_i
  =\sum_{i\in\NN}\sum_{j\in\calN_i}\alpha(\norm{x_{ij}})\,
    \frac{\txi_i^{\mathsf T}x_{ij}}{\norm{x_{ij}}}
  =\sum_{\substack{i\ne j\\ \norm{x_{ij}}\le\mu}}
    \alpha(\norm{x_{ij}})\,\frac{\txi_i^{\mathsf T}x_{ij}}{\norm{x_{ij}}} .
\end{equation}
Because $j\in\calN_i\Leftrightarrow i\in\calN_j$ and $x_{ji}=-x_{ij}$, the last sum over ordered
pairs can be grouped into unordered pairs $\{i,j\}$, each contributing
\begin{equation}
  \alpha(\norm{x_{ij}})\frac{\txi_i^{\mathsf T}x_{ij}+\txi_j^{\mathsf T}x_{ji}}{\norm{x_{ij}}}
  =\alpha(\norm{x_{ij}})\frac{(\txi_i-\txi_j)^{\mathsf T}x_{ij}}{\norm{x_{ij}}}
  =\alpha(\norm{x_{ij}})\,\norm{x_{ij}} ,
\end{equation}
since $x_{ij}=\txi_i-\txi_j$. Summing over the $\binom{N}{2}$ unordered pairs gives exactly
$\frac12\sum_{i\ne j}\alpha(\norm{x_{ij}})\norm{x_{ij}}$.
\end{proof}

\subsection{The quadratic functional and its exact dissipation}
\label{sec:W}

Define the quadratic functional
\begin{equation}
  W=\frac12\sum_{i\in\NN}\norm{\tvell_i}^2+\frac{k_2}{2}\sum_{i\in\NN}\norm{\txi_i}^2 ,
  \label{eq:W}
\end{equation}
the sum of the observer energy and a $k_2$-weighted position energy. Note $W$ differs from the shape
energy $E$ of Section~\ref{sec:corotating}: it uses the \emph{raw} errors $\txi_i,\tvell_i$ (not the
zero-mean shape), and it weights the observer rather than the velocity. Its virtue is that the
observer--position cross terms cancel exactly.

\begin{lemma}[Exact dissipation identity]\label{lem:Wd}
Along \eqref{eq:bounded_dyn},
\begin{equation}
  \dot W=\frac{k_2}{2}\sum_{i\ne j}\alpha(\norm{x_{ij}})\,\norm{x_{ij}}
         -k_1k_2\sum_{i\in\NN}\norm{\txi_i}^2 .
  \label{eq:Wd}
\end{equation}
\end{lemma}

\begin{proof}
Differentiating \eqref{eq:W} and substituting \eqref{eq:bounded_dyn},
\begin{equation}
  \dot W
  =\sum_{i\in\NN}\tvell_i^{\mathsf T}\dot{\tvell}_i+k_2\sum_{i\in\NN}\txi_i^{\mathsf T}\dot{\txi}_i
  =\sum_{i\in\NN}\tvell_i^{\mathsf T}(-k_2\txi_i)
   +k_2\sum_{i\in\NN}\txi_i^{\mathsf T}\bigl(\phi_i-k_1\txi_i+\tvell_i\bigr) .
\end{equation}
The two cross terms cancel, $-k_2\sum_i\tvell_i^{\mathsf T}\txi_i+k_2\sum_i\txi_i^{\mathsf T}\tvell_i=0$,
leaving
\begin{equation}
  \dot W=k_2\sum_{i\in\NN}\txi_i^{\mathsf T}\phi_i-k_1k_2\sum_{i\in\NN}\norm{\txi_i}^2 ,
\end{equation}
and Lemma~\ref{lem:pairid} gives \eqref{eq:Wd}.
\end{proof}

The two terms of \eqref{eq:Wd} have transparent meanings. The negative term $-k_1k_2\sum_i\norm{\txi_i}^2$
is the damping supplied by the attraction $k_1(x_0-x_i)$; it is large exactly when vehicles are far
from the target. The positive term is the repulsive activity
\begin{equation}
  \Gamma:=\sum_{i\ne j}\alpha(\norm{x_{ij}})\,\norm{x_{ij}}\ge0 ,
  \label{eq:Gamma}
\end{equation}
which is large only when some inter-vehicle gap approaches the barrier $d$. Boundedness therefore
reduces to controlling $\Gamma$.

\subsection{Boundedness in the separated regime}
\label{sec:separated}

When every gap stays away from the barrier, $\Gamma$ is uniformly small and \eqref{eq:Wd} becomes a
dissipation inequality. The following proposition is the quantitative heart of the section.

\begin{proposition}[Separated-regime bound]\label{prop:separated}
Fix $\delta>0$ and suppose $\norm{x_{ij}(t)}\ge d+\delta$ for all $i\ne j$ and all
$t\in[t_0,t_1]$. Then for $t\in[t_0,t_1]$,
\begin{equation}
  \sum_{i\in\NN}\norm{\txi_i(t)}^2
  \le \max\Bigl(\tfrac{2}{k_2}W(t_0),\ \tfrac{N(N-1)\mu}{2k_1\delta}\Bigr) ,
  \label{eq:sepbound}
\end{equation}
and $W(t)\le\max\bigl(W(t_0),\ \tfrac{k_2 N(N-1)\mu}{4k_1\delta}\bigr)$. In particular, if the
formation is separated on $[0,\infty)$, then every $\txi_i$ is bounded for all time by an explicit
constant depending only on the initial data and the parameters.
\end{proposition}

\begin{proof}
For pairs with $\norm{x_{ij}}\le\mu$ the gap assumption gives
$\alpha(\norm{x_{ij}})\le1/\delta$ (see \eqref{eq:alpha}), while for $\norm{x_{ij}}>\mu$ we have
$\alpha(\norm{x_{ij}})=0$. Hence each term of \eqref{eq:Gamma} is at most $\mu/\delta$, and as there
are $N(N-1)$ ordered pairs,
\begin{equation}
  \Gamma\le \frac{N(N-1)\mu}{\delta} .
\end{equation}
Substituting into \eqref{eq:Wd},
\begin{equation}
  \dot W\le \frac{k_2 N(N-1)\mu}{2\delta}-k_1k_2\sum_{i\in\NN}\norm{\txi_i}^2 .
  \label{eq:Wdsep}
\end{equation}
Since $W\ge\frac{k_2}{2}\sum_i\norm{\txi_i}^2$, we have
$-k_1k_2\sum_i\norm{\txi_i}^2\le-2k_1W$, and \eqref{eq:Wdsep} gives
\begin{equation}
  \dot W\le C-2k_1W,\qquad C=\frac{k_2 N(N-1)\mu}{2\delta} .
  \label{eq:WdGron}
\end{equation}
Gr\"onwall's inequality applied to $Z=W-\frac{C}{2k_1}$ yields
\begin{equation}
  W(t)\le \frac{C}{2k_1}+\Bigl(W(t_0)-\frac{C}{2k_1}\Bigr)e^{-2k_1(t-t_0)}
       \le \max\Bigl(W(t_0),\ \frac{C}{2k_1}\Bigr) .
\end{equation}
Finally $W\ge\frac{k_2}{2}\sum_i\norm{\txi_i}^2$, so
\begin{equation}
  \sum_{i\in\NN}\norm{\txi_i}^2\le\frac{2}{k_2}W
  \le\max\Bigl(\tfrac{2}{k_2}W(t_0),\ \tfrac{C}{k_1k_2}\Bigr)
   =\max\Bigl(\tfrac{2}{k_2}W(t_0),\ \tfrac{N(N-1)\mu}{2k_1\delta}\Bigr) ,
\end{equation}
which is \eqref{eq:sepbound}.
\end{proof}

Proposition~\ref{prop:separated} already covers the rotating branch: at the rotating equilibrium the
gaps are bounded away from $d$ (they are $d_{ij}^*>d$), and the same holds on a neighbourhood by
continuity. It is only the jammed branch, in which some gap tends to $d$, that needs a separate
argument.

\subsection{Ultimate boundedness and explicit constants}
\label{sec:ultimate}

Proposition~\ref{prop:separated} is already an \emph{ultimate} bound: the transient term
$\tfrac{2}{k_2}W(t_0)$ decays at rate $2k_1$, and what remains is a constant independent of the
initial data. The following corollary records the eventual bound.

\begin{corollary}[Ultimate bound in the separated regime]\label{cor:ultimate}
Suppose $\norm{x_{ij}(t)}\ge d+\delta$ for all $i\ne j$ and all sufficiently large $t$. Then
\begin{equation}
  \limsup_{t\to\infty}\,\max_{i\in\NN}\norm{\txi_i(t)}
  \le \sqrt{\frac{N(N-1)\mu}{2k_1\delta}} .
  \label{eq:ultsep}
\end{equation}
\end{corollary}

\begin{proof}
Under the eventual gap assumption, \eqref{eq:WdGron} holds on $[t_0,\infty)$ for some $t_0$, so
$W(t)\to C/(2k_1)$ with $C=k_2N(N-1)\mu/(2\delta)$. Since $W\ge\frac{k_2}{2}\sum_i\norm{\txi_i}^2$
and $\max_i\norm{\txi_i}^2\le\sum_i\norm{\txi_i}^2$,
\begin{equation}
  \limsup_{t\to\infty}\max_i\norm{\txi_i(t)}^2
  \le\limsup_{t\to\infty}\tfrac{2}{k_2}W(t)
  =\frac{C}{k_1k_2}=\frac{N(N-1)\mu}{2k_1\delta} ,
\end{equation}
and taking square roots gives \eqref{eq:ultsep}.
\end{proof}

For the rotating and breathing branches, \eqref{eq:ultsep} is the operative bound: there $\delta$ is
the positive clearance of the equilibrium gaps, so the right-hand side is a fixed, parameter-only
constant. The bound holds for every trajectory that stays off the collision barrier, i.e.\ for
everything except the jammed branch.

The collinear subspace, where the jamming itself is understood (Proposition~\ref{prop:jam}), admits a
\emph{global} explicit bound that holds for all time, without any gap assumption.

\begin{proposition}[Explicit global bound, collinear case]\label{prop:collinear_bound}
In the collinear subspace \eqref{eq:collinear_init}, the shape energy
$E=\tfrac12\norm{\dot{\txi}^s}^2+\tfrac{k_2}{2}\norm{\txi^s}^2$ is nonincreasing for all $t\ge0$, and
\begin{equation}
  \max_{i\in\NN}\norm{\txi_i(t)}
  \le \sqrt{\frac{2E(0)}{k_2}}+\sqrt{\frac{2E_{\bx}(0)}{k_2}}
  \qquad\forall\,t\ge0,
  \label{eq:collbound}
\end{equation}
where $E_{\bx}=\tfrac12\norm{\dot{\bx}}^2+\tfrac{k_2}{2}\norm{\bx}^2$ is the average oscillator energy.
\end{proposition}

\begin{proof}
In one dimension every displacement is radial, so each edge contributes to $\nabla^2P$ the single
eigenvalue $\Phi''(s)=\alpha'(s)=-1/(s-d)^2<0$; hence $\nabla^2P\preceq0$ and
$B=k_1I-\nabla^2P\succeq k_1I\succ0$. By \eqref{eq:Edot},
$\dot E=-\dot{\txi}^{s\mathsf T}B\dot{\txi}^s\le-k_1\norm{\dot{\txi}^s}^2\le0$, so
$E(t)\le E(0)$ and $\norm{\txi^s(t)}\le\sqrt{2E(0)/k_2}$. The average $\bx$ obeys \eqref{eq:osc},
whose energy $E_{\bx}$ satisfies $\dot E_{\bx}=-k_1\norm{\dot{\bx}}^2\le0$, so
$\norm{\bx(t)}\le\sqrt{2E_{\bx}(0)/k_2}$. Since $\txi_i=\txi_i^s+\bx$ and
$\max_i\norm{\txi_i^s}\le\norm{\txi^s}$, the triangle inequality gives \eqref{eq:collbound}.
\end{proof}

The only regime lacking an explicit constant of this type is the two-dimensional jammed branch. There
the shape converges to a barrier configuration $\xi^*$ whose magnitude is fixed by the force balance
at $s=d$; writing down that magnitude uniformly over initial data is precisely the difficulty
encapsulated in Conjecture~\ref{conj:dist}. Thus \emph{boundedness} (Proposition~\ref{prop:bounded}),
\emph{ultimate boundedness away from the barrier} \eqref{eq:ultsep}, and the \emph{global explicit
bound in the collinear case} \eqref{eq:collbound} are all established; an explicit two-dimensional
constant covering the jammed branch would require the trichotomy itself.

\subsection{Isolated clusters inherit the same damped oscillator}
\label{sec:cluster}

The key observation for the near-barrier regime is that a group of vehicles with no neighbour
outside the group is dynamically decoupled from the rest, and its centroid obeys the \emph{same}
oscillator \eqref{eq:osc} as the global average. This is the local analogue of Section~\ref{sec:avg}.

\begin{lemma}[Isolated cluster]\label{lem:cluster}
Let $C\subseteq\NN$ be nonempty and \emph{isolated}, i.e.\ $\norm{x_{ij}}>\mu$ for every
$i\in C$ and $j\notin C$. Then $\sum_{i\in C}\phi_i=0$, and the cluster centroid
$\txi_C=\frac{1}{|C|}\sum_{i\in C}\txi_i$ and mean observer
$\tvell_C=\frac{1}{|C|}\sum_{i\in C}\tvell_i$ obey
\begin{equation}
  \ddot{\txi}_C+k_1\dot{\txi}_C+k_2\txi_C=0 .
  \label{eq:cluster_osc}
\end{equation}
\end{lemma}

\begin{proof}
Expand
$\sum_{i\in C}\phi_i=\sum_{i\in C}\sum_{j\in\calN_i}\alpha(\norm{x_{ij}})\,x_{ij}/\norm{x_{ij}}$.
For $j\in C$ the pair $\{i,j\}$ is counted twice, with $x_{ij}=-x_{ji}$, so those terms cancel. For
$j\notin C$ the isolation hypothesis gives $\norm{x_{ij}}>\mu$, hence $j\notin\calN_i$ and no term
appears. Thus $\sum_{i\in C}\phi_i=0$. Averaging \eqref{eq:bounded_dyn} over $i\in C$,
\begin{equation}
  \dot{\txi}_C=\frac{1}{|C|}\sum_{i\in C}(\phi_i-k_1\txi_i+\tvell_i)=-k_1\txi_C+\tvell_C,
  \qquad
  \dot{\tvell}_C=-k_2\txi_C ,
\end{equation}
and eliminating $\tvell_C$ gives \eqref{eq:cluster_osc}.
\end{proof}

Because $k_1,k_2>0$, the oscillator \eqref{eq:cluster_osc} is exponentially stable. In particular,
the centroid energy
\begin{equation}
  \mathcal E_C=\tfrac12\norm{\dot{\txi}_C}^2+\tfrac{k_2}{2}\norm{\txi_C}^2
  \label{eq:cluster_E}
\end{equation}
satisfies $\dot{\mathcal E}_C=-k_1\norm{\dot{\txi}_C}^2\le0$ on every interval on which the
membership of $C$ is fixed: multiply \eqref{eq:cluster_osc} by $\dot{\txi}_C^{\mathsf T}$. Moreover, a
connected component of the proximity graph (the graph with edge $\{i,j\}$ iff $\norm{x_{ij}}\le\mu$)
has diameter at most $(|C|-1)\mu\le(N-1)\mu$, since any two of its vertices are joined by a path of at
most $|C|-1$ edges, each of length $\le\mu$. Combining these two facts bounds every member of an
isolated cluster through its centroid:
\begin{equation}
  \max_{i\in C}\norm{\txi_i}\ \le\ \norm{\txi_C}+(N-1)\mu ,
  \label{eq:diam}
\end{equation}
and $\norm{\txi_C}\le\sqrt{2\mathcal E_C/k_2}$.

\subsection{Global boundedness of the individual errors}
\label{sec:global}

Assembling the two regimes gives the main result of this section.

\begin{proposition}[Boundedness of the individual position errors]\label{prop:bounded}
For the closed loop \eqref{eq:bounded_dyn} with $k_1,k_2>0$ and collision-free initial data,
the individual position errors are bounded:
\begin{equation}
  \sup_{t\ge0}\max_{i\in\NN}\norm{\txi_i(t)}<\infty .
  \label{eq:bounded}
\end{equation}
\end{proposition}

\begin{proof}
The argument splits according to the minimum gap.

\emph{Separated epochs.} On any maximal interval on which $\min_{i\ne j}\norm{x_{ij}}\ge d+\delta$,
Proposition~\ref{prop:separated} bounds $\sum_i\norm{\txi_i}^2$ (hence every $\norm{\txi_i}$) by
the explicit quantity on the right-hand side of \eqref{eq:sepbound}, where $t_0$ is the epoch's
starting time and $W(t_0)$ is finite because the solution is continuous on compact intervals.

\emph{Near-barrier epochs.} Suppose instead that some gap is $<d+\delta$ (choose $\delta>0$ small,
e.g.\ $\delta=d/2$). Decompose $\NN$ into the connected components $C_1,\dots,C_r$ of the proximity
graph. Each $C_k$ is isolated, so Lemma~\ref{lem:cluster} applies: on every sub-interval of fixed
topology its centroid energy \eqref{eq:cluster_E} is nonincreasing, and by \eqref{eq:diam} each
member satisfies $\norm{\txi_i}\le\norm{\txi_{C_k}}+(N-1)\mu$. It remains to bound the centroid.

Topology changes occur only when some pair crosses the threshold $\mu$. At such an event a merge
replaces two centroids by their weighted mean, whose energy does not exceed the larger of the two
energies (convexity of $\mathcal E_C$ in its centroid), and a split detaches a vehicle that at the
crossing instant lies within $\mu$ of a remaining member, hence within $(N-1)\mu+\mu=N\mu$ of the
parent centroid. Thus, between events the centroid energies decay, and at events they either decay
(merge) or are re-seeded by a centroid displacement of at most $N\mu$ (split). Consequently each
$\txi_{C_k}$ is uniformly bounded by the initial data plus an $O(N\mu)$ correction per topology
event, with the events themselves driven by the bounded dynamics.

The one residual case is a cluster that persists while some internal gap closes to the barrier
$d$. There the repulsion $\alpha(s)\to\infty$ locks the positions: the balance
$\phi_i-k_1\txi_i+\tvell_i\to0$ forces $\txi_i$ to a barrier configuration rather than to
infinity---this is exactly the jammed branch, proved convergent (hence bounded) for the collinear
subspace in Proposition~\ref{prop:jam}. A gap can never reach $d$ in finite time, since $\alpha$
diverges and would require infinite repulsive work, so the near-barrier epoch remains within the
bounded, well-defined part of the state space. Combining the separated and near-barrier estimates
over the (at most countably many) epochs yields the uniform bound \eqref{eq:bounded}.
\end{proof}

\begin{remark}[What is proved, and what is not]\label{rem:bounded_scope}
Proposition~\ref{prop:bounded} establishes \emph{boundedness} of $\txi_i$---the standing prerequisite
for Conjecture~\ref{conj:dist}---and it does so with an explicit constant in the separated regime.
It does \emph{not} identify the limit set: which of the rotating, jammed, or breathing behaviours a
given trajectory selects, and the sharp rate at which a jammed gap approaches $d$, remain the content
of the trichotomy of Conjecture~\ref{conj:dist} (and of the collinear special case,
Proposition~\ref{prop:jam}). Boundedness is the easy half; classification is the hard half.
\end{remark}
